\begin{abstract}
Automated cervical-cytology screening is increasingly framed as an image-classification problem, and deep convolutional networks now report near-ceiling accuracy on curated single-dataset benchmarks. Whether such accuracy reflects transferable diagnostic competence or merely fits the idiosyncrasies of one slide-preparation protocol remains largely unaudited. We present a rigorous, leakage-aware audit of a knowledge-distilled cervical-cytology classifier. An EfficientNet-B4 teacher (18.5\,M parameters) was trained on the five SIPaKMeD cytomorphology classes with weighted cross-entropy and two-stage fine-tuning, and its knowledge was distilled into an EfficientNet-B0 student (4.67\,M parameters) using the Hinton objective \citep{hinton2015distilling}. In-distribution, the distilled student surpasses its teacher on every metric, reaching accuracy 0.966, balanced accuracy 0.968, macro-F1 0.967, Cohen's $\kappa$ 0.958, one-vs-rest macro AUROC 0.998, and an Expected Calibration Error of 0.015; Grad-CAM \citep{selvaraju2017gradcam} confirms that decisions localize to the nucleus and cell body rather than background. We then evaluate the same student, without retraining, on three external Papanicolaou datasets---Herlev, APACC, and HiCervix---under a deterministic, documented binary Normal-versus-Abnormal label alignment with reported excluded fractions. Cross-dataset accuracy collapses to 0.639 (Herlev), 0.605 (HiCervix), and 0.532 (APACC), falling below the majority-class baseline on two of three datasets despite all AUROC exceeding chance. A diagnostic analysis attributes this collapse not to label-mapping artefacts---class-conditional $P(\text{abnormal})$ separation is positive everywhere and every AUROC $>0.5$---but to operating-point miscalibration that does not transfer across domains and to increasing slide-preparation distance from SIPaKMeD's conventional smear. We further show that the internal accuracy is itself optimistic, inflated by same-field-of-view leakage in the per-cell split. The study contributes a reproducible cross-dataset audit methodology, evidence that distillation can yield a smaller-yet-stronger in-distribution model that nonetheless fails out-of-distribution, and concrete, evidence-led remedies for trustworthy cervical-screening deployment.
\end{abstract}

\section{Introduction}

Cervical cancer remains one of the most preventable malignancies, and its prevention hinges on the cytological screening of Papanicolaou-stained smears. Because expert cytologists are scarce and manual reading is laborious and subject to inter-observer variability, automated classification of cervical cells has become an active research frontier. Modern deep convolutional networks routinely report accuracies above 0.95 on public benchmarks such as SIPaKMeD \citep{plissiti2018sipakmed}, and this success has fuelled optimism that screening can be partially automated. Yet the dominant evaluation paradigm in this literature is single-dataset: a model is trained and tested on splits drawn from one collection, prepared under one staining and segmentation protocol, and the resulting accuracy is reported as if it were an intrinsic property of the model. This conflates two very different questions---how well a model fits the statistics of a particular slide-preparation pipeline, and how well it would actually generalize to the heterogeneous slides encountered in clinical practice. The gap between these questions is precisely where deployable, trustworthy screening tools succeed or fail, and it is rarely measured.

This paper subjects a representative high-accuracy system to exactly such a measurement. We study a knowledge-distilled classifier in which an EfficientNet-B4 teacher \citep{tan2019efficientnet}, ImageNet-pretrained \citep{deng2009imagenet} and trained on the five SIPaKMeD cytomorphology classes, is distilled into a compact EfficientNet-B0 student through the temperature-scaled Hinton objective \citep{hinton2015distilling}. The student is attractive for deployment because it carries only 4.67\,M parameters against the teacher's 18.5\,M, and we ask whether such compression sacrifices accuracy, calibration, explainability, or generalization. We evaluate the model both within SIPaKMeD and, crucially, across three independent external Papanicolaou datasets---Herlev \citep{jantzen2005herlev}, APACC, and HiCervix \citep{cai2024hicervix}---which differ in staining, segmentation, and slide preparation (conventional smear, auto-segmented crops, and liquid-based cytology). Because the five SIPaKMeD morphology classes do not share a label space with the dysplasia and Bethesda \citep{nayar2015bethesda} schemes of the external datasets, we align them at the only common, deterministically mappable level: a binary Normal-versus-Abnormal decision, with glandular and out-of-scope categories explicitly excluded and the excluded fraction reported as an honesty metric.

Our investigation is organized around four research questions and their hypotheses. \textbf{RQ1}: Can a knowledge-distilled B0 student match the B4 teacher's in-distribution accuracy at a fraction of the parameters? \textbf{H1}: yes, and the regularizing effect of distillation may let the student match or even exceed the teacher. \textbf{RQ2}: Does strong in-distribution accuracy translate to cross-dataset generalization on independent Pap datasets? \textbf{H2}: no---a large internal-to-external gap is expected. \textbf{RQ3}: Is any cross-dataset failure a label-mapping artefact or a genuine model and domain failure? \textbf{H3}: it is genuine domain shift compounded by miscalibration, not a labelling error. \textbf{RQ4}: Do calibration and explainability hold, and does calibration transfer out-of-distribution? \textbf{H4}: in-distribution the model is well-calibrated \citep{guo2017calibration} and its Grad-CAM attributions are plausible, but calibration does not transfer across domains. Together these questions move the evaluation beyond a single accuracy figure toward an account of where, and why, a high-performing model breaks.

In answering them, this work makes the following contributions.

\begin{itemize}
\item A reproducible, leakage-aware cross-dataset audit of a distilled cervical-cytology model across three external Papanicolaou datasets (Herlev, APACC, HiCervix) under a deterministic, documented binary label alignment with reported excluded fractions.
\item A diagnostic methodology that separates labelling artefacts from genuine domain failure, combining class-conditional $P(\text{abnormal})$ separation, an AUROC-above-chance test, and an accuracy-versus-majority-baseline and oracle-threshold analysis.
\item Empirical evidence that the distilled student surpasses its teacher while remaining well-calibrated and explainable in-distribution, yet collapses out-of-distribution, with the collapse attributed to operating-point miscalibration and slide-preparation distance rather than to relabelling \citep{geirhos2020shortcut}.
\item Concrete, evidence-led remedies: leakage-safe cluster-grouped evaluation, stain and domain normalization \citep{macenko2009stain}, per-domain calibration, and the reporting of threshold-free AUROC alongside the excluded fraction.
\end{itemize}